% answers_short.tex —— 2007 年·上海卷·理科 · 简短答案速查
% 由 tools/make_short_answers.py 自 answers.tex 的 \gksolhead/\gkanswer 提取，
% 并与 2007/2007_数学真题_上海理_答案要点.md 逐题比对后生成，勿手改（改源 MD 后重跑）。

\begin{gknotice}
\textbf{使用说明}：本节只给最终答案便于快速核对，完整推导见解析版。答案为编写者独立推导，\textbf{非官方标准答案}。
\end{gknotice}

\gkpart{填空题}
\begin{itemize}[leftmargin=2.6em]
  \item[\textbf{1.}] $(-\infty,3)\cup(3,4)$。
  \item[\textbf{2.}] $-\frac23$。
  \item[\textbf{3.}] $f^{-1}(x)=\frac{x}{x-1},\ x\ne1$。
  \item[\textbf{4.}] $x=\log_3 7$。
  \item[\textbf{5.}] $\frac1{16}$。
  \item[\textbf{6.}] $\pi$。
  \item[\textbf{7.}] $\frac3{10}$。
  \item[\textbf{8.}] $y^2=12(x+3)$。
  \item[\textbf{9.}] ②④。
  \item[\textbf{10.}] 例如 $s_1\parallel s_2$ 且 $t_1,t_2$ 相交。
  \item[\textbf{11.}] $d=2\sin\theta,\ 0<\theta<\pi$。
\end{itemize}

\gkpart{选择题}
\begin{center}\begin{tabular}{ccccc}
\toprule
题号 & 12 & 13 & 14 & 15 \\
\midrule
答案 & A & C & B & D \\
\bottomrule
\end{tabular}\end{center}

\gkpart{解答题}
\begin{itemize}[leftmargin=2.6em]
  \item[\textbf{16.}] $\arcsin\frac1{\sqrt6}$。
  \item[\textbf{17.}] $\frac87$。
  \item[\textbf{18.}] （1）$2499.8$ 兆瓦；（2）至少 $61.5\%$。
  \item[\textbf{19.}] （1）$a=0$ 时为偶函数，其余非奇非偶；（2）$a\le16$。
  \item[\textbf{20.}] （1）$2,5,8,11,8,5,2$；（2）$k=13$，最大 $626$；（3）记 $W=(1,2,\ldots,2^{m-1})$，四种为 $W$ 后接去首项的 $W^R$、$W^R$ 后接去首项的 $W$、$WW^R$、$W^RW$。选择第一种，$m\ge2008$ 时 $S_{2008}=2^{2008}-1$；$1500<m<2008$ 时 $S_{2008}=3\cdot2^{m-1}-2^{2m-2009}-1$。
  \item[\textbf{21.}] （1）$a^2=7/4,b^2=1,c^2=3/4$；（2）$1<b/c<4/3$；（3）仅 $k=0$。
\end{itemize}
